\documentclass[11pt]{article}
% Bien dich: pdflatex (hoac Overleaf). Moi hinh ve viet bang TikZ/pgfplots ngay trong file, khong can anh ngoai.
\usepackage[T5]{fontenc}
\usepackage[utf8]{inputenc}
\usepackage[vietnamese]{babel}
\usepackage[a4paper,margin=2.3cm]{geometry}
\usepackage{graphicx} % Required for inserting images
\usepackage{amsmath,amssymb}
\usepackage{booktabs}
\usepackage{array}
\usepackage{float}
\usepackage{caption}
\usepackage{xcolor}
\usepackage{tikz}
\usetikzlibrary{positioning,arrows.meta,shapes.geometric,fit,calc}
\usepackage{pgfplots}
\pgfplotsset{compat=1.17}
\usepackage[hidelinks]{hyperref}

\definecolor{cOrig}{HTML}{7F7F7F}
\definecolor{cV1}{HTML}{1F77B4}
\definecolor{cV2}{HTML}{D62728}
\definecolor{cSS}{HTML}{FF7F0E}
\definecolor{cBox}{HTML}{EAF2FB}
\definecolor{cNew}{HTML}{FDECEA}

\newcommand{\R}{\mathbb{R}}
\newcommand{\bmu}{\boldsymbol{\mu}}
\newcommand{\bSigma}{\boldsymbol{\Sigma}}
\newcommand{\vf}{\mathbf{f}}
\newcommand{\tr}{\operatorname{tr}}

\title{Báo cáo dự án khoa học\\[4pt]
\large Học tăng dần theo lớp không lưu mẫu cũ trên ViT tiền huấn luyện:\\
từ RSIAT đến V1 (AGME) và V2 (AGME + STS)}
\author{Thị Thanh Nga Hoàng}
\date{November 2026}

\begin{document}

\maketitle

\begin{abstract}
Báo cáo trình bày pipeline hiện tại của dự án học tăng dần theo lớp (class-incremental learning, CIL) không lưu mẫu
cũ, xây trên ViT-B/16 tiền huấn luyện với adapter. Xuất phát từ phương pháp nền RSIAT, chúng tôi phân tích sai số của
bộ phân loại dựng từ thống kê lớp thành hai phần: phần do \emph{biểu diễn} và phần do \emph{ước lượng thống kê} của lớp
cũ. Từ phân tích này, chúng tôi đề xuất hai phiên bản. \textbf{V1 (AGME)} đơn giản hoá RSIAT (bỏ $\mathcal{L}_{RS}$ và
$\mathcal{L}_{orth}$) và thay bộ phân loại đơn bằng tổ hợp chuẩn hoá z-score của ba bộ phân loại có sai số khác nhau
(CA, LDA và AGM -- một hồi quy ridge dạng đóng trên đặc trưng ngẫu nhiên, tính chính xác từ mô-men Gaussian).
\textbf{V2} giữ nguyên V1 nhưng thay bù drift SDC bằng \textbf{STS} (Self-Tracking Statistics): ước lượng
responsibility bằng bước E của mô hình trộn có thành phần nền, dò lớp thiếu dữ liệu bằng mẫu đối kháng, khớp ánh xạ
drift (Procrustes, chọn độ co bằng split-half) và \emph{vận chuyển cả mean lẫn covariance} của lớp cũ. Trên
ImageNet-A (3 seed), V2 đạt $73.62/65.94$ (trung bình/task cuối) so với $72.68/65.02$ của RSIAT gốc
($+0.94/+0.92$); trên VTAB đạt $94.99/94.91$ so với $94.72/94.69$. Trên CUB, riêng bộ phân loại CA của V2 vượt RSIAT
gốc ($94.16/90.54$ so với $93.93/90.37$), nhưng tổ hợp ba bộ phân loại làm giảm điểm -- chúng tôi phân tích nguyên
nhân và giới hạn của phương pháp.
\end{abstract}

\tableofcontents
\newpage

% =====================================================================================================================
\section{Giới thiệu}

\paragraph{Bài toán.} Trong học tăng dần theo lớp, mô hình nhận lần lượt các \emph{task}; mỗi task mang một nhóm lớp
mới chưa từng thấy. Sau mỗi task, mô hình phải phân loại được \emph{mọi} lớp đã học, không được biết ảnh thuộc task
nào. Ràng buộc \emph{không lưu mẫu cũ} (exemplar-free) cấm giữ lại ảnh của các task trước -- một ràng buộc thực tế vì
lý do bộ nhớ và quyền riêng tư. Khó khăn trung tâm là \textbf{quên thảm khốc}: khi học lớp mới, biểu diễn thay đổi
(\emph{drift}) và kiến thức về lớp cũ bị phá vỡ.

\paragraph{Vì sao dùng mô hình tiền huấn luyện.} Các nghiên cứu gần đây cho thấy bắt đầu từ ViT tiền huấn luyện
(ImageNet-21k) và chỉ tinh chỉnh một phần nhỏ tham số (adapter, prompt) giúp giảm drift đáng kể
\cite{slca,ssiat}. Khi đó, lớp cũ có thể được ``nhớ'' qua \emph{thống kê đặc trưng} -- mean và covariance của từng lớp
-- thay vì qua ảnh; bộ phân loại được huấn luyện lại trên mẫu sinh từ các Gaussian này (classifier alignment, CA
\cite{slca}).

\paragraph{Câu hỏi của dự án.} Nếu lớp cũ chỉ còn tồn tại dưới dạng thống kê, thì độ chính xác cuối cùng bị giới hạn
bởi (i) biểu diễn còn tách được các lớp tốt đến đâu, và (ii) thống kê đã lưu còn đúng với biểu diễn hiện tại đến đâu.
Dự án đi từ phân tích này để cải tiến phương pháp nền \textbf{RSIAT}, qua hai phiên bản:
\begin{itemize}
  \item \textbf{V1 (AGME v1)}: đơn giản hoá RSIAT và thay bộ phân loại đơn bằng tổ hợp ba bộ phân loại dựng từ cùng
        thống kê (mục~\ref{sec:v1}).
  \item \textbf{V2}: V1 với cơ chế bù drift mới STS, ước lượng lại cả mean lẫn covariance của lớp cũ
        (mục~\ref{sec:v2}).
\end{itemize}

\paragraph{Đóng góp của báo cáo.}
(1) Mô tả đầy đủ pipeline, từ vấn đề, công thức tới lý do của từng thành phần.
(2) Phân rã sai số làm cơ sở thiết kế, kèm bằng chứng thực nghiệm (thí nghiệm ``oracle'').
(3) Kết quả đầy đủ đến hết V2 trên ImageNet-A, VTAB, CUB, cùng các thí nghiệm loại bỏ (ablation) và các hướng đã thử
nhưng thất bại.

% =====================================================================================================================
\section{Thiết lập bài toán và ký hiệu}\label{sec:setup}

Chuỗi task $\mathcal{T}_0,\dots,\mathcal{T}_{T-1}$; task $t$ có tập lớp $\mathcal{C}_t$ (không giao nhau) và dữ liệu
$\mathcal{D}_t=\{(x_i,y_i)\}$, $y_i\in\mathcal{C}_t$. Ký hiệu $K_t=|\mathcal{C}_0\cup\dots\cup\mathcal{C}_{t-1}|$ là
số lớp cũ khi bắt đầu task $t$. Bộ trích đặc trưng $f_\theta:\mathcal{X}\to\R^{d}$ ($d=768$) là ViT-B/16 tiền huấn
luyện ImageNet-21k, đóng băng, với adapter có thể học ở mỗi khối (\texttt{ffn\_num}$=64$). Ký hiệu $f_{t}$ là bộ trích
đặc trưng sau task $t$.

Mỗi lớp $c$ đã học được nhớ bằng thống kê $(\bmu_c,\bSigma_c)$, tính trên ảnh huấn luyện (biến đổi kiểu test) qua mô
hình tại thời điểm lớp được học:
\begin{equation}
  \bmu_c=\frac{1}{n_c}\sum_{i:y_i=c} f(x_i),\qquad
  \bSigma_c=\frac{1}{n_c-1}\sum_{i:y_i=c}\big(f(x_i)-\bmu_c\big)\big(f(x_i)-\bmu_c\big)^{\!\top}+10^{-3}I .
\end{equation}
Khi biểu diễn thay đổi ở các task sau, $(\bmu_c,\bSigma_c)$ cần được \emph{bù drift} để khớp không gian đặc trưng mới.

\paragraph{Thước đo.} Sau mỗi task $t$, đo độ chính xác $A_t$ trên tập test của mọi lớp đã học. Báo cáo
\textbf{Avg} $=\frac1T\sum_t A_t$ (trung bình qua các task) và \textbf{Last} $=A_{T-1}$ (sau task cuối).

% =====================================================================================================================
\section{Phương pháp nền: RSIAT}\label{sec:rsiat}

RSIAT thuộc họ phương pháp adapter-tuning kèm bù drift ngữ nghĩa (gần với SSIAT \cite{ssiat}). Mỗi task gồm bốn bước.

\paragraph{(a) Huấn luyện adapter với CosFace.} Head cosine với trọng số $w_j$; với $\cos\theta_j=\langle
\hat{\vf},\hat w_j\rangle$ (dấu mũ là chuẩn hoá $\ell_2$), hàm mất mát CosFace \cite{cosface} có lề $m$ và thang $s$:
\begin{equation}
  \mathcal{L}_{\text{CE}}=-\log\frac{e^{s(\cos\theta_y-m)}}{e^{s(\cos\theta_y-m)}+\sum_{j\neq y}e^{s\cos\theta_j}},
\end{equation}
chỉ tính trên các lớp của task hiện tại. Lề $m$ làm cụm lớp chặt hơn, giúp thống kê Gaussian về sau sát thực tế hơn.

\paragraph{(b) Ràng buộc chống quên.}
\begin{itemize}
\item $\mathcal{L}_{RS}$ (chỉ task 0): trên đặc trưng chuẩn hoá $\hat{\vf}$, kéo các cặp cùng lớp về tích vô hướng $1$
và đẩy các cặp khác lớp dưới lề $\rho$:
$\mathcal{L}_{RS}=\overline{[1-\hat\vf_i^\top\hat\vf_j]_+}^{\,y_i=y_j}+\alpha\,\overline{[\hat\vf_i^\top\hat\vf_j-\rho]_+}^{\,y_i\neq y_j}$.
\item \textbf{RAE} (residual autoencoder), từ task 1: đặc trưng mới phải gần đặc trưng cũ \emph{sau một phép biến
đổi học được}:
\begin{equation}
  \mathcal{L}_{RAE}=\beta\,\big\|f_t(x)-g_\phi\big(f_{t-1}(x)\big)\big\|_2^2,\qquad
  g_\phi(u)=u+\sigma\big(\mathrm{Dec}(\mathrm{Enc}(u))\big),
\end{equation}
với Enc: $768\to64\to$ \texttt{code\_dims}, Dec ngược lại, $\sigma$ là sigmoid. $g_\phi$ được học cùng adapter. So với
chưng cất đặc trưng cứng ($g=\mathrm{id}$), RAE cho phép một độ lệch dương có cấu trúc, nên đặc trưng ``được neo''
nhưng không bị đóng băng.
\item $\mathcal{L}_{orth}$: phạt độ tương đồng cosine trung bình giữa prototype lớp cũ (qua $g_\phi$) và đặc trưng cũ
(qua $g_\phi$), với trọng số $\gamma$.
\end{itemize}
Tổng mất mát tại task $t\ge1$: $\mathcal{L}=\mathcal{L}_{CE}+\mathcal{L}_{RAE}+\gamma\mathcal{L}_{orth}$; tại task 0:
$\mathcal{L}_{CE}+\lambda_{RS}\mathcal{L}_{RS}$.

\paragraph{(c) Bù drift SDC.} Semantic Drift Compensation \cite{sdc} ước lượng dịch chuyển của mean lớp cũ từ dịch
chuyển của ảnh \emph{mới}: với $\Delta_i=f_t(x_i)-f_{t-1}(x_i)$ trên ảnh của task $t$,
\begin{equation}
  \bmu_c\leftarrow\bmu_c+\sum_i w_{ic}\Delta_i,\qquad
  w_{ic}\propto\exp\!\Big(-\frac{\|f_{t-1}(x_i)-\bmu_c\|^2}{2\sigma^2}\Big),\ \ \sigma=4 .
\end{equation}
SDC \emph{không cập nhật covariance}: $\bSigma_c$ giữ nguyên như lúc lớp được học.

\paragraph{(d) Classifier alignment (CA).} Sau khi có thống kê mọi lớp, đóng băng backbone và huấn luyện lại head
cosine $10$ epoch trên mẫu giả: mỗi epoch, mỗi lớp $c$ sinh $256$ mẫu
$z\sim\mathcal{N}\big((0.9+\delta_c)\bmu_c,\ \bSigma_c\big)$, với $\delta_c=0.1\,(\tau_c+1)/(t+1)$ và $\tau_c$ là task
học lớp $c$. Hệ số $<1$ cho lớp cũ là một heuristic chống thiên lệch về lớp mới (với head cosine, co mean tương đương
nới rộng phân phối theo hướng).

% =====================================================================================================================
\section{Cơ sở phân tích: sai số của bộ phân loại dựng từ thống kê}\label{sec:theory}

Mọi bộ phân loại trong pipeline (CA, LDA, AGM ở dưới) đều được dựng từ thống kê lớp. Với một bộ phân loại
``plug-in'' dựng từ phân phối ước lượng, sai số tách được thành hai phần \cite[Định lý~2.2]{devroye}:
\begin{equation}
  \underbrace{\mathrm{Err}}_{\text{sai số cuối}} \;=\;
  \underbrace{\mathrm{Err}^{\star}(f_t)}_{\text{(I) do biểu diễn}} \;+\;
  \underbrace{\big(\mathrm{Err}-\mathrm{Err}^{\star}(f_t)\big)}_{\text{(II) do ước lượng thống kê sai}},
\end{equation}
trong đó $\mathrm{Err}^\star(f_t)$ là sai số tốt nhất có thể với biểu diễn hiện tại, và phần (II) bị chặn bởi độ lệch
giữa phân phối ước lượng và phân phối thật \cite{devroye}.

\begin{itemize}
\item \textbf{Phần (I)} chỉ phụ thuộc vào việc các lớp cũ còn tách khỏi nhau và khỏi lớp mới tốt đến đâu. Ràng buộc
chống quên (RAE) kiểm soát phần này.
\item \textbf{Phần (II)} phụ thuộc vào việc $(\bmu_c,\bSigma_c)$ còn đúng không sau drift. Bù drift (SDC, STS) và cách
dùng thống kê (bộ phân loại, tổ hợp) kiểm soát phần này.
\end{itemize}
Thí nghiệm \emph{oracle} đo trực tiếp phần (II): thay thống kê lớp cũ bằng thống kê thật (tính lại từ ảnh train cũ qua
mô hình hiện tại -- chỉ dùng để phân tích, không phải phương pháp). Khoảng cách oracle -- thực tế là phần (II) còn có
thể lấy lại. Phân rã này dẫn dắt toàn bộ thiết kế: V1 tác động vào \emph{cách dùng} thống kê, V2 tác động vào
\emph{độ đúng} của thống kê, còn RAE được giữ để không làm phần (I) xấu đi (mục~\ref{sec:history} cho thấy bỏ RAE làm
phần (I) tăng mạnh).

% =====================================================================================================================
\section{V1 (AGME v1): đơn giản hoá và tổ hợp ba bộ phân loại}\label{sec:v1}

\subsection{Đơn giản hoá backbone}
V1 dùng \textbf{RSIAT-simple}: bỏ $\mathcal{L}_{RS}$ ($\lambda_{RS}=0$) và $\mathcal{L}_{orth}$ ($\gamma=0$), giữ
nguyên CosFace, RAE (cùng $\beta$), SDC và CA. Lý do: hai thành phần này làm pipeline phức tạp và có thêm siêu tham số
riêng cho từng dataset; khi bỏ đi, độ chính xác CA gần như không đổi (ImageNet-A seed 1993: CA $74.76/65.64$ so với
$74.76/65.77$ của RSIAT gốc; CUB: $93.99/90.46$ so với $93.93/90.37$), nên phần cải thiện về sau đến từ các thành phần
mới chứ không từ việc tune lại.

\subsection{Ba bộ phân loại dùng chung thống kê}
Cho đặc trưng test $\vf$ và thống kê $\{(\bmu_c,\bSigma_c)\}_{c=1}^{C}$:
\begin{enumerate}
\item \textbf{CA}: logit của head cosine sau bước CA (mục~\ref{sec:rsiat}d).
\item \textbf{LDA} \cite{esl}: covariance chung có ridge,
\begin{equation}
 \bar\bSigma=\frac1C\sum_c\bSigma_c,\quad P=\big(\bar\bSigma+\overline{\mathrm{diag}}(\bar\bSigma)\,I\big)^{-1},\quad
 \ell^{LDA}_c(\vf)=\vf^\top P\bmu_c-\tfrac12\bmu_c^\top P\bmu_c .
\end{equation}
Ridge bằng trung bình đường chéo là co mạnh -- phù hợp chế độ $d=768$ lớn so với số ảnh mỗi lớp ($\approx30$--$50$).
\item \textbf{AGM} (analytic Gaussian-moment head): hồi quy ridge trên đặc trưng ngẫu nhiên
\begin{equation}
 \phi(\vf)=\big[\vf;\ \mathrm{ReLU}(W\vf);\ 10\big]\in\R^{M},\qquad W\in\R^{4000\times d},\ W_{ij}\sim\mathcal{N}(0,1/d)
 \ \text{(cố định)},
\end{equation}
($M=4000+768+1$), nghiệm dạng đóng
\begin{equation}
 G=\sum_c \bar n\,\mathbb{E}_{c}[\phi\phi^\top],\quad q_c=\bar n\,\mathbb{E}_{c}[\phi],\quad
 w_c=(G+\lambda I)^{-1}q_c,\quad \lambda=0.1\,\tr(G)/M,
\end{equation}
trong đó kỳ vọng lấy theo $\vf\sim\mathcal{N}(\bmu_c,\bSigma_c)$ và được tính \emph{chính xác, không lấy mẫu}: khối
ReLU--ReLU dùng mô-men của phân phối chuẩn hai biến bị cắt (cầu phương Gauss--Legendre), khối tuyến tính--ReLU dùng
bổ đề Stein; $\bar n$ là số ảnh trung bình mỗi lớp (trọng số cân bằng). Đây là phiên bản ``vô hạn mẫu'' của việc huấn
luyện bộ phân loại trên mẫu Gaussian, với đặc trưng phi tuyến kiểu random features \cite{rahimi}.
\end{enumerate}

\subsection{Tổ hợp z-score}
Ba bộ phân loại có thang logit khác nhau, nên mỗi vector logit được chuẩn hoá theo từng ảnh rồi cộng với trọng số bằng
nhau:
\begin{equation}
 z(\ell)=\frac{\ell-\mathrm{mean}_c(\ell)}{\mathrm{std}_c(\ell)},\qquad
 \hat y=\arg\max_c\ \big[z(\ell^{CA})+z(\ell^{LDA})+z(\ell^{AGM})\big]_c .
\end{equation}
\textbf{Lý do:} CA là bộ phân loại phân biệt (học từ mẫu), LDA là bộ sinh tuyến tính với covariance chung, AGM là bộ
phi tuyến dùng covariance riêng từng lớp; sai số của chúng tương quan không hoàn toàn nên tổ hợp giảm phương sai (đặc
biệt ở lớp cũ, nơi thống kê kém tin cậy). Không có trọng số nào được chọn trên tập test.

% =====================================================================================================================
\section{V2: STS -- thống kê tự bám drift}\label{sec:v2}

\subsection{Động cơ}
SDC có hai hạn chế: (i) trọng số $w_{ic}$ là kernel khoảng cách cố định, nên ảnh mới không liên quan (ở xa mọi lớp cũ)
vẫn ``kéo'' mean lớp cũ; (ii) covariance không được cập nhật, trong khi CA, LDA, AGM đều dùng covariance. Thí nghiệm
oracle trên ImageNet-A (mục~\ref{sec:oracle}) cho thấy phần (II) không nhỏ: thay thống kê lớp cũ bằng thống kê thật
làm tổ hợp CA+LDA tăng từ $74.91/66.10$ lên $75.53/67.08$. STS (Self-Tracking Statistics) được thiết kế để thu hẹp
khoảng cách này.

V2 $=$ V1 với SDC được thay bằng STS (chạy \emph{một lần ở cuối mỗi task}). RAE và mọi thành phần khác giữ nguyên. Vì
STS không tham gia hàm mất mát, quỹ đạo huấn luyện backbone của V1 và V2 là \emph{như nhau} (mọi tính toán của STS bảo
toàn trạng thái RNG); khác biệt kết quả đến hoàn toàn từ thống kê.

\begin{figure}[t]
\centering
\begin{tikzpicture}[
  font=\small, node distance=6mm and 8mm,
  box/.style={draw, rounded corners=2pt, fill=cBox, align=center, minimum height=9mm, inner sep=4pt},
  nbox/.style={box, fill=cNew},
  arr/.style={-{Stealth[length=2.2mm]}, thick}]
\node[box] (data) {Ảnh train\\task $t$};
\node[box, right=of data] (train) {Huấn luyện adapter\\$\mathcal{L}_{CE}+\beta\mathcal{L}_{RAE}$};
\node[box, above=of train] (old) {Mô hình cũ $f_{t-1}$\\(đóng băng)};
\node[nbox, right=of train] (sts) {\textbf{STS} (V2), dùng $f_{t-1}$ và $f_t$\\bước E + mẫu đối kháng\\ánh xạ drift\\vận chuyển $\bmu,\bSigma$};
\node[box, above=of sts] (bank) {Kho thống kê\\$\{(\bmu_c,\bSigma_c)\}$};
\node[box, right=of sts] (newst) {Thống kê lớp mới\\(qua $f_t$)};
\node[box, below=10mm of sts] (ca) {CA: head cosine\\trên mẫu Gaussian};
\node[nbox, left=of ca] (lda) {LDA};
\node[nbox, right=of ca] (agm) {AGM};
\node[nbox, below=of ca] (fuse) {$z(\ell^{CA})+z(\ell^{LDA})+z(\ell^{AGM})$ (V1, V2)};
\draw[arr] (data) -- (train);
\draw[arr] (old) -- (train);
\draw[arr] (train) -- (sts);
\draw[arr] (bank) -- (sts);
\draw[arr] (sts) -- (newst);
\draw[arr] (newst.north) |- (bank.east);
\draw[arr] (bank.west) -| ($(lda.north)+(0,0)$);
\draw[arr] (sts) -- (ca);
\draw[arr] (newst.south) |- (agm.east);
\draw[arr] (lda) |- (fuse.west);
\draw[arr] (ca) -- (fuse);
\draw[arr] (agm) |- (fuse.east);
\end{tikzpicture}
\caption{Pipeline cho một task $t\ge1$. Ô màu đỏ nhạt là phần thêm so với RSIAT: STS (V2) thay SDC; LDA, AGM và tổ hợp
z-score (V1). Ô xanh giữ như RSIAT-simple.}
\label{fig:pipeline}
\end{figure}

\subsection{Tập dò và bước E với thành phần nền}
Tại đầu task $t$, đặt $F=\{f_{t-1}(x_i)\}$ là vị trí của ảnh train mới qua mô hình \emph{cũ}. Các lớp cũ được mô hình
hoá bằng mô hình trộn Gaussian, cộng một \emph{thành phần nền} cho chính ảnh mới:
\begin{equation}
 p(\vf)=\sum_{c=1}^{K}\frac{1-\pi_0}{K}\,\mathcal{N}(\vf;\bmu_c,\bSigma_{pool})+\pi_0\,\mathcal{N}(\vf;\bmu_0,\bSigma_0),
\end{equation}
với $\bSigma_{pool}$ là trung bình $\bSigma_c$ cộng ridge như LDA, $(\bmu_0,\bSigma_0)$ khớp trên ảnh mới, $\pi_0=0.5$.
Bước E \cite{em} cho \emph{responsibility} $r_{ic}=p(c\mid \vf_i)$ và khối lượng $N_c=\sum_i r_{ic}$.

\textbf{Lý do:} khác SDC, ảnh mới không giống lớp cũ nào sẽ rơi vào thành phần nền và không làm lệch ước lượng; trọng
số do chính mô hình thống kê quyết định thay vì một bán kính kernel cố định.

\paragraph{Mẫu đối kháng (cho lớp thiếu dữ liệu).} Lớp có $N_c<5$ gần như không có ảnh mới nào ``đại diện''. Theo
tinh thần ADC \cite{adc}, chọn $16$ ảnh mới gần $\bmu_c$ nhất và chạy $10$ bước dấu-gradient (bước $2/255$) qua mô hình
cũ để giảm khoảng cách Mahalanobis tới $\bmu_c$; chỉ giữ mẫu mà bước E gán vào đúng lớp $c$. Các mẫu này chỉ dùng để
\emph{đo} drift, không đưa vào hàm mất mát.

\subsection{Ánh xạ drift}
Cuối task, $F'=\{f_t(x_i)\}$ là vị trí của cùng tập dò qua mô hình \emph{mới}, $D=F'-F$.
\begin{enumerate}
\item \textbf{Bước E bám:} $r\leftarrow(1-\eta)\,r+\eta\,r^{new}$ với $\eta=0.3$.
\item \textbf{Ánh xạ chung (Procrustes):} khớp $D_i\approx m_D+(\vf_i-\bar m)B_g$ với $I+B_g$ là phép quay
\cite{procrustes} có co về $I$:
\begin{equation}
 I+B_g=\arg\max_{R^\top R=I}\tr\!\Big(R^\top\big(X_c^\top Y_c+\alpha\tfrac{\|X_c\|_F^2}{d}I\big)\Big),
\end{equation}
($X_c$, $Y_c$ là $F$ và $F'-\bar m_D$ đã trừ trung bình). $\alpha\in\{10^{-2},\dots,10^{2}\}$ chọn bằng
\emph{split-half}: khớp trên một nửa ảnh của mỗi lớp mới, chọn $\alpha$ có sai số dự đoán nhỏ nhất trên nửa còn lại.
\textbf{Lý do:} phép quay giữ hình dạng covariance, ít bậc tự do (tránh khớp quá) -- phù hợp drift nhỏ khi có RAE.
\item \textbf{Dịch chuyển riêng từng lớp:} với lớp có $N_c\ge5$, ước lượng tịnh tiến cục bộ có trọng số $r_{ic}$:
$d^{loc}_c=\bar D_c+(\bmu_c-\bar x_c)B_g$, rồi co về dự đoán chung bằng $n_0=50$ ``mẫu giả'':
\begin{equation}
 d_c=\frac{n_{\text{eff},c}\,d^{loc}_c+n_0\,d^{glob}_c}{n_{\text{eff},c}+n_0},\qquad
 n_{\text{eff},c}=\frac{(\sum_i r_{ic})^2}{\sum_i r_{ic}^2}.
\end{equation}
Đây là ước lượng co kiểu Bayes phân cấp: lớp có nhiều bằng chứng dùng dự đoán riêng, lớp ít bằng chứng dựa vào ánh xạ
chung.
\end{enumerate}

\subsection{Vận chuyển mean và covariance}
\begin{equation}
 \bmu_c\leftarrow\bmu_c+d_c,\qquad \bSigma_c\leftarrow(I+B_g)^\top\,\bSigma_c\,(I+B_g).
\end{equation}
\textbf{Lý do:} nếu $\vf'=\vf(I+B)+b$ thì $\mathrm{Cov}(\vf')=(I+B)^\top\mathrm{Cov}(\vf)(I+B)$ -- covariance được
biến đổi nhất quán với mean, điều SDC không làm. Hình~\ref{fig:transport} minh hoạ.

\begin{figure}[t]
\centering
\begin{tikzpicture}[font=\small]
\draw[thick, cOrig, rotate around={20:(0,0)}] (0,0) ellipse (1.4 and 0.6);
\fill[cOrig] (0,0) circle (1.5pt) node[below=2pt] {$\bmu_c$};
\node[cOrig] at (0,1.2) {không gian $f_{t-1}$};
\draw[thick, cV2, rotate around={55:(5,0.6)}] (5,0.6) ellipse (1.4 and 0.6);
\fill[cV2] (5,0.6) circle (1.5pt) node[below right=1pt] {$\bmu_c+d_c$};
\node[cV2] at (5,2.0) {không gian $f_{t}$ (STS)};
\draw[dashed, thick, cV1, rotate around={20:(4.2,0.1)}] (4.2,0.1) ellipse (1.4 and 0.6);
\node[cV1, align=center] at (3.3,-1.3) {SDC: chỉ dời mean,\\covariance giữ nguyên};
\draw[-{Stealth[length=2.5mm]}, thick] (0.3,0.1) to[bend left=15] node[above] {$\vf\mapsto \vf(I+B_g)+b$} (4.7,0.6);
\foreach \p in {(-0.6,-0.2),(0.5,0.3),(0.9,0.5),(-0.3,0.1),(0.2,-0.3)} \fill[black!50] \p circle (1pt);
\foreach \p in {(4.7,0.0),(5.3,1.2),(5.6,1.3),(4.9,0.5),(5.1,0.1)} \fill[black!50] \p circle (1pt);
\end{tikzpicture}
\caption{STS vận chuyển cả mean và covariance qua cùng ánh xạ drift; SDC (nét đứt) chỉ dời mean. Chấm xám: ảnh mới
được dùng làm tập dò.}
\label{fig:transport}
\end{figure}

\subsection{Thuật toán và các lựa chọn thiết kế}
\noindent\fbox{\begin{minipage}{0.96\linewidth}
\textbf{Thuật toán V2 tại task $t\ge1$}
\begin{enumerate}\itemsep0pt
\item \emph{Đầu task:} tính $F$ qua $f_{t-1}$; bước E $\Rightarrow r, N_c$; sinh mẫu đối kháng cho lớp có $N_c<5$.
\item \emph{Huấn luyện:} $\mathcal{L}_{CE}+\beta\mathcal{L}_{RAE}$ (như V1), STS không can thiệp.
\item \emph{Cuối task:} tính $F'$ qua $f_t$; bước E bám; ánh xạ Procrustes ($\alpha$ split-half); dịch chuyển riêng
      có co; vận chuyển $(\bmu_c,\bSigma_c)$ của mọi lớp cũ.
\item \emph{Lớp mới:} $(\bmu_c,\bSigma_c)$ qua $f_t$.
\item \emph{CA} trên mẫu Gaussian của mọi lớp; \emph{dự đoán} bằng $z(\ell^{CA})+z(\ell^{LDA})+z(\ell^{AGM})$.
\end{enumerate}
\end{minipage}}

\medskip
Các lựa chọn đã được kiểm chứng bằng thực nghiệm (mục~\ref{sec:history}):
\begin{itemize}
\item \textbf{Giữ RAE.} Bỏ RAE và để STS ``bám'' drift lớn (cùng một hàm phạt phần drift STS không giải thích được)
chỉ hoà với V1; bỏ RAE hoàn toàn mất $2.5/3.6$ điểm. Theo phân rã ở mục~\ref{sec:theory}, RAE giữ phần (I), STS giảm
phần (II): hai vai trò bổ sung, không thay thế nhau.
\item \textbf{Bám một lần cuối task} thay vì sau mỗi epoch: cho cùng kết quả ($75.63/66.69$ so với $75.67/66.69$) mà
nhanh hơn khoảng $3$ lần; STS khi đó là một mô-đun thay SDC, độc lập với vòng huấn luyện.
\end{itemize}

% =====================================================================================================================
\section{Thiết lập thực nghiệm}\label{sec:exp}

\paragraph{Dữ liệu và giao thức.}
\textbf{ImageNet-A} (200 lớp, 10 task $\times$ 20 lớp, $\approx30$ ảnh train/lớp);
\textbf{VTAB} (50 lớp, 5 task $\times$ 10 lớp, mỗi task một miền, $\approx36$ ảnh/lớp);
\textbf{CUB-200} (200 lớp, 20 task $\times$ 10 lớp, $\approx47$ ảnh/lớp). Thứ tự lớp xáo theo seed (ImageNet-A, CUB);
VTAB giữ thứ tự miền cố định. Backbone ViT-B/16 ImageNet-21k, adapter, SGD với lịch cosine.

\begin{table}[H]
\centering\small
\caption{Siêu tham số theo dataset (giống RSIAT gốc; V1/V2 đặt $\lambda_{RS}=\gamma=0$). STS dùng chung:
$\eta=0.3$, $\pi_0=0.5$, ngưỡng $N_c=5$, $n_0=50$, $16$ mẫu đối kháng/lớp, $10$ bước, ánh xạ Procrustes,
$\alpha\in\{10^{-2},\dots,10^2\}$.}
\begin{tabular}{lccc}
\toprule
 & ImageNet-A & VTAB & CUB \\
\midrule
Lớp/task, số task & 20, 10 & 10, 5 & 10, 20 \\
Epoch task 0 / task sau / CA & 20 / 15 / 10 & 30 / 20 / 10 & 20 / 30 / 10 \\
Batch, lr & 48, $6.7\!\times\!10^{-3}$ & 48, $3.6\!\times\!10^{-2}$ & 128, $3.8\!\times\!10^{-2}$ \\
CosFace $s$ / $m$ & 30 / 0.3 & 20 / 0.3 & 20 / 0.1 \\
RAE $\beta$ / \texttt{code\_dims} & 3.0 / 128 & 0.2 / 256 & 1.5 / 768 \\
RSIAT gốc: $\lambda_{RS}$ / $\gamma$ & 1.0 / 6.0 & 2.0 / 1.0 & 0.2 / 0.75 \\
\bottomrule
\end{tabular}
\end{table}

\paragraph{Ghi chú về độ tin cậy.}
\begin{itemize}
\item Huấn luyện là tất định (cùng cấu hình cho kết quả lặp lại chính xác). Trên ImageNet-A, tập test nhỏ: một ảnh
tương đương $0.06$--$0.3$ điểm tuỳ task, nên chênh lệch dưới $\approx0.3$ trên một seed nằm trong mức nhiễu.
\item \textbf{VTAB:} seed chỉ dùng để xáo thứ tự lớp, còn seed của torch cố định ($=1$); VTAB không xáo (thứ tự miền
cố định), nên ``các seed'' của VTAB cho \emph{cùng một} lần chạy. Các con số ``3 seed'' của VTAB trước đây khác nhau
chỉ vì mã lúc đó chưa hoàn toàn tất định. Vì vậy VTAB chỉ báo cáo một lần chạy.
\item Số RSIAT gốc: ImageNet-A và VTAB là chạy lại trong cùng môi trường; CUB là số do tác giả cung cấp (log RSIAT gốc
trong kho mã không khớp với cấu hình gốc). RSIAT gốc chỉ có bộ phân loại CA.
\item Con số V1/V2 được tính offline trên dump (\texttt{exps/agm/final\_eval.py}); cột CA tính từ dump có thể lệch log
huấn luyện $\le0.1$ điểm.
\end{itemize}

% =====================================================================================================================
\section{Kết quả}\label{sec:results}

\subsection{ImageNet-A}

\begin{table}[H]
\centering\small
\caption{ImageNet-A, seed 1993: độ chính xác sau mỗi task. RSIAT gốc: CA; V1, V2: tổ hợp z(CA)+z(LDA)+z(AGM).
In đậm: cao nhất mỗi cột.}
\label{tab:ina1993}
\setlength{\tabcolsep}{4pt}
\begin{tabular}{l*{10}{c}cc}
\toprule
Task & 0 & 1 & 2 & 3 & 4 & 5 & 6 & 7 & 8 & 9 & Avg & Last \\
\midrule
RSIAT gốc & 86.29 & 82.50 & 79.62 & 78.65 & 75.51 & 73.23 & 70.56 & \textbf{68.59} & 66.88 & 65.77 & 74.76 & 65.77 \\
RSIAT-super simple (bỏ RAE) & 87.43 & 83.33 & 78.57 & 76.07 & 74.08 & 71.28 & 66.64 & 65.06 & 63.94 & 62.34 & 72.88 & 62.34 \\
V1 & \textbf{87.43} & \textbf{84.72} & \textbf{81.09} & \textbf{79.08} & 75.03 & 74.36 & 71.01 & 68.51 & 66.81 & 65.90 & 75.39 & 65.90 \\
V2 & \textbf{87.43} & 84.44 & 80.88 & 78.51 & \textbf{75.62} & \textbf{74.87} & \textbf{71.63} & \textbf{68.99} & \textbf{67.24} & \textbf{66.69} & \textbf{75.63} & \textbf{66.69} \\
\bottomrule
\end{tabular}
\end{table}

\begin{figure}[H]
\centering
\begin{tikzpicture}
\begin{axis}[width=0.92\linewidth, height=6.2cm, xlabel={Task}, ylabel={Độ chính xác (\%)},
  xtick={0,...,9}, ymin=61, ymax=89, grid=major, grid style={black!12},
  legend style={at={(0.02,0.03)}, anchor=south west, font=\small, fill=white, fill opacity=0.9}]
\addplot[cOrig, thick, mark=o] coordinates {(0,86.29)(1,82.50)(2,79.62)(3,78.65)(4,75.51)(5,73.23)(6,70.56)(7,68.59)(8,66.88)(9,65.77)};
\addlegendentry{RSIAT gốc (CA)}
\addplot[cSS, thick, dashed, mark=triangle] coordinates {(0,87.43)(1,83.33)(2,78.57)(3,76.07)(4,74.08)(5,71.28)(6,66.64)(7,65.06)(8,63.94)(9,62.34)};
\addlegendentry{Super simple (bỏ RAE)}
\addplot[cV1, thick, mark=square] coordinates {(0,87.43)(1,84.72)(2,81.09)(3,79.08)(4,75.03)(5,74.36)(6,71.01)(7,68.51)(8,66.81)(9,65.90)};
\addlegendentry{V1}
\addplot[cV2, very thick, mark=*] coordinates {(0,87.43)(1,84.44)(2,80.88)(3,78.51)(4,75.62)(5,74.87)(6,71.63)(7,68.99)(8,67.24)(9,66.69)};
\addlegendentry{V2}
\end{axis}
\end{tikzpicture}
\caption{ImageNet-A, seed 1993. V2 bắt đầu vượt V1 từ task 4 (khi drift tích luỹ đủ lớn để việc vận chuyển covariance
có tác dụng) và hơn RSIAT gốc $+0.92$ ở task cuối. Bỏ RAE (cam) làm mất $3.6$ điểm ở task cuối.}
\label{fig:ina1993}
\end{figure}

\begin{table}[H]
\centering\small
\caption{ImageNet-A, 3 seed (Avg / Last). V2 thắng RSIAT gốc ở cả 3 seed.}
\label{tab:ina3}
\begin{tabular}{lcccc}
\toprule
Seed & RSIAT gốc & V1 & V2 & V2 $-$ gốc \\
\midrule
1993 & 74.76 / 65.77 & 75.39 / 65.90 & \textbf{75.63 / 66.69} & $+0.87$ / $+0.92$ \\
1994 & 72.18 / 63.99 & 72.98 / 64.98 & \textbf{73.06 / 65.04} & $+0.88$ / $+1.05$ \\
1995 & 71.10 / 65.31 & \textbf{72.19} / 65.70 & 72.18 / \textbf{66.10} & $+1.08$ / $+0.79$ \\
\midrule
\textbf{Trung bình} & 72.68 / 65.02 & 73.52 / 65.53 & \textbf{73.62 / 65.94} & $\mathbf{+0.94}$ / $\mathbf{+0.92}$ \\
\bottomrule
\end{tabular}
\end{table}

\begin{figure}[H]
\centering
\begin{tikzpicture}
\begin{axis}[ybar, bar width=12pt, width=0.8\linewidth, height=5.6cm, ymin=64, ymax=74.5,
  symbolic x coords={Avg,Last}, xtick=data, enlarge x limits=0.5, ylabel={Độ chính xác (\%)},
  nodes near coords, nodes near coords style={font=\scriptsize, rotate=90, anchor=west},
  legend style={at={(0.5,1.03)}, anchor=south, legend columns=3, font=\small}]
\addplot[fill=cOrig!70] coordinates {(Avg,72.68)(Last,65.02)};
\addplot[fill=cV1!70] coordinates {(Avg,73.52)(Last,65.53)};
\addplot[fill=cV2!80] coordinates {(Avg,73.62)(Last,65.94)};
\legend{RSIAT gốc, V1, V2}
\end{axis}
\end{tikzpicture}
\caption{ImageNet-A, trung bình 3 seed.}
\end{figure}

\begin{table}[H]
\centering\small
\caption{ImageNet-A, seed 1993: từng bộ phân loại (Avg / Last). Từ V1 sang V2, cả ba bộ phân loại đều tăng ở task
cuối, nhiều nhất là LDA ($+0.66$) và AGM ($+0.85$) -- hai bộ dùng covariance mà SDC không cập nhật.}
\label{tab:heads}
\begin{tabular}{lccccc}
\toprule
 & CA & LDA & AGM & z(CA+AGM) & z(CA+LDA+AGM) \\
\midrule
V1 & 74.76 / 65.64 & 74.33 / 64.98 & 74.19 / 64.65 & 75.25 / 65.90 & 75.39 / 65.90 \\
V2 & 75.02 / 66.23 & 74.42 / 65.64 & 74.55 / 65.50 & 75.50 / 66.29 & \textbf{75.63 / 66.69} \\
\bottomrule
\end{tabular}
\end{table}

\paragraph{Nhận xét.} (i) Tổ hợp luôn tốt hơn bộ phân loại tốt nhất đơn lẻ ($+0.3$ ở V1, $+0.5$ ở V2 tại task cuối). (ii) Ở
task 1--3, V2 nhỉnh thua V1 ($-0.2$ đến $-0.6$) dù CA bằng hoặc cao hơn; từ task 4 trở đi V2 thắng ở mọi task.
(iii) Phần đóng góp của STS so với V1 là $+0.10/+0.41$ trung bình 3 seed: chủ yếu ở task cuối.

\subsection{VTAB}

\begin{table}[H]
\centering\small
\caption{VTAB (một lần chạy tất định). RSIAT gốc: CA; V1, V2: tổ hợp.}
\begin{tabular}{l*{5}{c}cc}
\toprule
Task & 0 & 1 & 2 & 3 & 4 & Avg & Last \\
\midrule
RSIAT gốc & 96.70 & 93.93 & 94.00 & \textbf{94.27} & 94.69 & 94.72 & 94.69 \\
V1 & \textbf{97.10} & 94.29 & 94.17 & 94.14 & 94.86 & 94.91 & 94.86 \\
V2 & \textbf{97.10} & \textbf{94.36} & \textbf{94.34} & 94.25 & \textbf{94.91} & \textbf{94.99} & \textbf{94.91} \\
\bottomrule
\end{tabular}
\end{table}

\begin{table}[H]
\centering\small
\caption{VTAB: từng bộ phân loại của V2 và hai biến thể tăng khả năng học miền mới (Avg / Last).}
\begin{tabular}{lccccc}
\toprule
 & CA & LDA & AGM & z(CA+AGM) & z(CA+LDA+AGM) \\
\midrule
V2 & 94.80 / 94.73 & 93.99 / 93.94 & 94.34 / 94.42 & 94.97 / 94.89 & \textbf{94.99 / 94.91} \\
V2, 30 epoch/task & 94.93 / 94.64 & 93.62 / 93.42 & 94.12 / 94.05 & 94.92 / 94.65 & 94.93 / 94.67 \\
V2, $\beta=0$ (bỏ RAE) & 94.61 / 94.20 & 94.17 / 94.25 & 94.12 / 94.16 & 94.77 / 94.87 & 94.78 / 94.85 \\
\bottomrule
\end{tabular}
\end{table}

\paragraph{Lỗi còn lại nằm ở đâu.} Ở task cuối của V2 (sai $\approx5.3\%$): nhầm giữa các miền $\approx0.7\%$, quên
(độ chính xác lúc vừa học trừ lúc cuối) $\approx0.5\%$, \textbf{nhầm giữa các lớp trong cùng miền $\approx4.6\%$}; miền
của task 1 chỉ đạt $86.75\%$ ngay khi vừa học. Như vậy trên VTAB, dư địa nằm ở chất lượng học từng miền (phần (I)),
không ở thống kê lớp cũ; điều này giải thích vì sao STS chỉ thêm $+0.08/+0.05$ so với V1. Bỏ RAE làm CA giảm
($94.20$ so với $94.73$ ở task cuối) -- RAE không cản việc học miền mới.

\subsection{CUB-200}

\begin{table}[H]
\centering\small
\caption{CUB, seed 1993 (Avg / Last). RSIAT gốc do tác giả cung cấp.}
\label{tab:cub}
\begin{tabular}{lccccc}
\toprule
 & CA & LDA & AGM & z(CA+AGM) & z(CA+LDA+AGM) \\
\midrule
RSIAT gốc & 93.93 / 90.37 & -- & -- & -- & -- \\
V1 & 93.99 / 90.46 & 93.16 / 88.72 & 93.40 / 89.06 & 93.63 / 89.40 & 93.66 / 89.57 \\
V2 & \textbf{94.16 / 90.54} & 93.38 / 89.10 & 93.50 / 89.65 & 93.67 / 89.99 & 93.69 / 90.08 \\
\bottomrule
\end{tabular}
\end{table}

\begin{figure}[H]
\centering
\begin{tikzpicture}
\begin{axis}[width=0.92\linewidth, height=6cm, xlabel={Task}, ylabel={Độ chính xác (\%)},
  xtick={0,2,...,19}, ymin=89.3, ymax=100.3, grid=major, grid style={black!12},
  legend style={at={(0.98,0.97)}, anchor=north east, font=\small}]
\addplot[cOrig, thick, mark=o, mark size=1.5pt] coordinates {(0,100.00)(1,98.79)(2,98.55)(3,98.02)(4,97.50)(5,95.24)(6,93.70)(7,94.23)(8,93.82)(9,93.27)(10,93.40)(11,92.58)(12,91.79)(13,91.65)(14,91.23)(15,91.43)(16,91.21)(17,91.02)(18,90.80)(19,90.37)};
\addlegendentry{RSIAT gốc (CA)}
\addplot[cV2, very thick, mark=*, mark size=1.5pt] coordinates {(0,100.00)(1,98.79)(2,98.55)(3,98.02)(4,97.68)(5,95.39)(6,94.09)(7,94.34)(8,94.20)(9,93.60)(10,93.40)(11,92.94)(12,92.19)(13,92.20)(14,91.68)(15,91.80)(16,91.36)(17,91.35)(18,91.02)(19,90.54)};
\addlegendentry{V2 (chỉ CA)}
\addplot[cV2, thick, dashed, mark=x, mark size=2pt] coordinates {(0,99.15)(1,98.38)(2,98.26)(3,97.80)(4,97.68)(5,95.39)(6,93.83)(7,94.12)(8,93.92)(9,93.27)(10,93.25)(11,92.37)(12,91.33)(13,91.29)(14,90.95)(15,90.96)(16,90.77)(17,90.70)(18,90.35)(19,90.08)};
\addlegendentry{V2 (tổ hợp)}
\addplot[cV1, thick, dashed, mark=square, mark size=1.3pt] coordinates {(0,99.15)(1,98.38)(2,98.26)(3,98.02)(4,98.04)(5,95.39)(6,93.83)(7,94.34)(8,94.11)(9,93.27)(10,93.09)(11,92.44)(12,91.14)(13,91.29)(14,90.89)(15,90.96)(16,90.57)(17,90.37)(18,90.13)(19,89.57)};
\addlegendentry{V1 (tổ hợp)}
\end{axis}
\end{tikzpicture}
\caption{CUB, seed 1993. CA của V2 bằng hoặc cao hơn RSIAT gốc ở cả 20 task; tổ hợp (nét đứt) thấp hơn CA.}
\label{fig:cub}
\end{figure}

\paragraph{Phân tích.} (i) STS có ích: chỉ khác nhau ở bù drift, CA của V2 hơn CA của V1 $+0.17$ trung bình, đều ở
mọi task (lớn nhất $+0.61$ ở task 13); LDA và AGM của V2 hơn V1 $+0.38/+0.59$ ở task cuối. (ii) Tổ hợp làm giảm điểm
ở \emph{cả} V1 và V2: trên CUB (phân loại chim rất mịn), LDA và AGM yếu hơn CA rõ rệt ($1.4$--$1.7$ điểm ở task cuối
của V1; ngay ở task 0, khi chưa có drift, LDA chỉ đạt $98.31$ so với $100$ của CA). Cộng đều ba bộ phân loại vì thế kéo
về phía bộ yếu. Đây là giới hạn của quy tắc tổ hợp trọng số bằng nhau, không phải của STS.

\subsection{Thí nghiệm oracle: phần (II) còn lại bao nhiêu}\label{sec:oracle}
Trên ImageNet-A seed 1993 (V2), so ba loại thống kê lớp cũ trên \emph{cùng} backbone, với tổ hợp z(CA)+z(LDA):

\begin{table}[H]
\centering\small
\begin{tabular}{lccc}
\toprule
Thống kê lớp cũ & CA & z(CA+LDA) & Sai số mean lớp cũ (chỉ số tương đối) \\
\midrule
Không bù (giữ nguyên lúc học) & 74.93 / 65.96 & 74.91 / 66.10 & 4.90 / 6.27 \\
STS (V2) & 75.02 / 66.23 & 75.13 / 65.77 & 2.55 / 3.43 \\
Oracle (thống kê thật) & 75.10 / 66.23 & 75.53 / 67.08 & 0 / 0 \\
\bottomrule
\end{tabular}
\end{table}
STS giảm sai số mean lớp cũ khoảng một nửa. Khoảng cách tới oracle ($\approx0.4$ trung bình, $\approx1.3$ ở task
cuối với z(CA+LDA)) là dư địa còn lại của phần (II); phần lớn sai số còn lại thuộc phần (I).

\subsection{Lịch sử phát triển và các hướng đã thử}\label{sec:history}

\paragraph{Bỏ RAE.} RSIAT-super simple (RSIAT-simple bỏ RAE, không cơ chế giữ nào): $72.88/62.34$ so với $75.39/65.90$
của V1, tức $-2.51/-3.56$. Với CA ở task cuối, độ chính xác lớp cũ giảm từ $65.45$ xuống $61.36$ trong khi lớp mới tăng từ
$68.55$ lên $77.42$: bỏ ràng buộc thì học lớp mới dễ hơn nhưng quên lớp cũ nhanh, khoảng cách tăng dần theo task
(Bảng~\ref{tab:ina1993}, Hình~\ref{fig:ina1993}).

\paragraph{STS thay RAE (``bám thay vì neo'').} Ý tưởng: bỏ RAE, cho backbone tự do, STS bám drift sau mỗi epoch, và
chỉ phạt phần drift mà STS không giải thích được: $\beta\,\mathrm{pen}\big(f_t(x),T(f_{t-1}(x))\big)$ với $T$ là ánh
xạ tích luỹ. Kết quả trên ImageNet-A seed 1993:

\begin{table}[H]
\centering\small
\caption{Các biến thể STS không có RAE (z(CA+LDA+AGM), Avg / Last; các run chạy đủ 10 task).}
\begin{tabular}{lcc}
\toprule
Biến thể & Avg / Last & So với V1 \\
\midrule
V1 (tham chiếu) & 75.39 / 65.90 & -- \\
Affine, MSE, $\beta=3$ & 75.37 / 65.96 & $-0.02$ / $+0.06$ \\
Affine, MSE, $\beta=3$, bám 3 lần/epoch & 75.34 / 65.70 & $-0.05$ / $-0.20$ \\
Procrustes + metric Fisher, $\beta=3$ & 75.46 / 65.90 & $+0.07$ / $0.00$ \\
RAE + STS mỗi epoch & 75.67 / 66.69 & $+0.28$ / $+0.79$ \\
\textbf{RAE + STS cuối task (= V2)} & \textbf{75.63 / 66.69} & $+0.24$ / $+0.79$ \\
\bottomrule
\end{tabular}
\end{table}

Các biến thể dừng sớm vì kém rõ rệt (CA tại task 3, V1: $78.37$): $\beta=1$ affine $77.36$; Procrustes $\beta=1$ và
$\beta=3$ $77.79$; Fisher $\beta=10$ $75.79$; Fisher $\beta=30$ \emph{phân kỳ}; neo bằng mẫu đối kháng làm lớp mới
giảm $\approx2.6$ điểm; replay Gaussian lớp cũ trong lúc train làm task 1 giảm $2.5$ điểm; neo Fisher (đích là đặc
trưng cũ) thua V1 từ task 2. Thí nghiệm replay oracle cho thấy thống kê STS bám được gần thống kê thật, nhưng
\emph{biểu diễn} lớp cũ đã xấu đi -- đúng như phân rã ở mục~\ref{sec:theory}: bỏ RAE làm phần (I) tăng nhiều hơn phần
(II) giảm. Kết luận dẫn tới V2: giữ RAE cho phần (I), dùng STS cho phần (II).

% =====================================================================================================================
\section{Thảo luận}\label{sec:discussion}

\paragraph{Điều đã làm được.} V2 vượt RSIAT gốc ổn định trên ImageNet-A (3/3 seed, $+0.94/+0.92$) và VTAB
($+0.27/+0.22$), và trên CUB nếu chỉ dùng CA ($+0.23/+0.17$). Mỗi thành phần có lý do rõ: tổ hợp ba bộ phân loại khai
thác sai số khác nhau; STS sửa hai điểm yếu của SDC (trọng số không phân biệt ảnh liên quan, covariance không được cập
nhật).

\paragraph{Giới hạn.}
\begin{itemize}
\item Mức cải thiện nhỏ ($<1$ điểm); trên VTAB và CUB mới có một lần chạy.
\item Tổ hợp trọng số bằng nhau làm hại trên CUB; cần cách gán trọng số theo độ tin cậy của từng bộ phân loại, chọn
      trên dữ liệu train giữ lại (không dùng test).
\item Phân rã sai số cho thấy phần (I) chiếm ưu thế: các cải tiến về thống kê đang tiến gần trần (khoảng cách oracle
      trên ImageNet-A chỉ còn $\approx0.4$ trung bình).
\item Một chi tiết mã cần ghi nhận: ở task 0, learning rate của các nhóm tham số được đặt cố định $0.01$ (giá trị
      \texttt{init\_lr} trong cấu hình không có tác dụng ở task 0). Mọi phiên bản chịu như nhau nên so sánh vẫn công bằng.
\end{itemize}

\paragraph{Vị trí so với các công trình liên quan.} Từng thành phần của STS có tiền lệ: ước lượng drift từ ảnh mới
(SDC \cite{sdc}), mẫu đối kháng để đo drift (ADC \cite{adc}), ánh xạ đặc trưng cũ--mới (LDC \cite{ldc}). Điểm riêng của
STS là bước E với thành phần nền, dịch chuyển riêng từng lớp co về ánh xạ chung theo kiểu Bayes phân cấp với độ co chọn
bằng split-half, và vận chuyển nhất quán cả mean lẫn covariance. Để khẳng định đóng góp, cần so trực tiếp với ADC và LDC
trên cùng backbone.

% =====================================================================================================================
\section{Kết luận}
Báo cáo trình bày pipeline từ RSIAT tới V2. Dựa trên phân rã sai số thành phần do biểu diễn và phần do ước lượng thống
kê, V1 cải thiện cách \emph{dùng} thống kê (tổ hợp CA, LDA, AGM), V2 cải thiện \emph{độ đúng} của thống kê lớp cũ (STS).
V2 đạt $73.62/65.94$ trên ImageNet-A (3 seed), $94.99/94.91$ trên VTAB, và vượt RSIAT gốc trên CUB ở bộ phân loại CA.
Các bước tiếp theo: (i) gán trọng số tổ hợp theo độ tin cậy cho CUB, (ii) chạy đủ 3 seed và thêm CIFAR-100, ImageNet-R,
OmniBenchmark, (iii) so sánh với ADC và LDC. Hướng V3 (thống kê có độ bất định) đang được phát triển riêng và không
thuộc phạm vi báo cáo này.

% =====================================================================================================================
\begin{thebibliography}{99}
\bibitem{slca} G. Zhang, L. Wang, G. Kang, L. Chen, Y. Wei. SLCA: Slow Learner with Classifier Alignment for Continual
Learning on a Pre-trained Model. \emph{ICCV}, 2023.
\bibitem{ssiat} Y. Tan, Q. Zhou, X. Xiang, K. Wang, Y. Wu, Y. Li. Semantically-Shifted Incremental Adapter-Tuning is A
Continual ViTransformer. \emph{CVPR}, 2024.
\bibitem{sdc} L. Yu, B. Twardowski, X. Liu, L. Herranz, K. Wang, Y. Cheng, S. Jui, J. van de Weijer. Semantic Drift
Compensation for Class-Incremental Learning. \emph{CVPR}, 2020.
\bibitem{adc} D. Goswami, A. Soutif-Cormerais, Y. Liu, S. Kamath, B. Twardowski, J. van de Weijer. Resurrecting Old
Classes with New Data for Exemplar-Free Continual Learning. \emph{CVPR}, 2024.
\bibitem{ldc} A. Gomez-Villa, D. Goswami, K. Wang, A. D. Bagdanov, B. Twardowski, J. van de Weijer. Exemplar-free
Continual Representation Learning via Learnable Drift Compensation. \emph{ECCV}, 2024.
\bibitem{cosface} H. Wang, Y. Wang, Z. Zhou, X. Ji, D. Gong, J. Zhou, Z. Li, W. Liu. CosFace: Large Margin Cosine Loss
for Deep Face Recognition. \emph{CVPR}, 2018.
\bibitem{devroye} L. Devroye, L. Györfi, G. Lugosi. \emph{A Probabilistic Theory of Pattern Recognition}. Springer, 1996.
\bibitem{esl} T. Hastie, R. Tibshirani, J. Friedman. \emph{The Elements of Statistical Learning}, 2nd ed., Springer, 2009.
\bibitem{rahimi} A. Rahimi, B. Recht. Random Features for Large-Scale Kernel Machines. \emph{NeurIPS}, 2007.
\bibitem{em} A. P. Dempster, N. M. Laird, D. B. Rubin. Maximum Likelihood from Incomplete Data via the EM Algorithm.
\emph{J. Royal Statistical Society B}, 1977.
\bibitem{procrustes} P. H. Schönemann. A Generalized Solution of the Orthogonal Procrustes Problem. \emph{Psychometrika},
1966.
\end{thebibliography}

\end{document}
